\newpage
\section{Fase 3 - PatchAgent: il ciclo iterativo e la FSM}
\label{sec:3.4}
\label{sec:fase3}
La terza e ultima fase dell'architettura \texttt{HybridRepair},
implementata dal modulo \texttt{PatchAgent}, costituisce il nucleo
esecutivo del sistema.

Le tre fasi si ripartiscono tre domande distinte. La Fase 1 (\texttt{FaultOracle})
localizza il difetto e ne individua le cause, rispondendo a
\textit{dove} si trova l'errore e \textit{perché} il programma fallisce.
La Fase 2 (\texttt{SpecReason}) astrae queste informazioni in una specifica
concettuale di riparazione, rispondendo a \textit{cosa} debba essere
modificato. Il \texttt{PatchAgent} cala questi costrutti nel codice sorgente,
rispondendo a \textit{come} il codice Java vada concretamente alterato.

Questa traduzione da specifica concettuale a modifica
concreta non è banale. Il \texttt{PatchAgent}
deve infatti garantire che la modifica generata
soddisfi tre criteri rigorosi:
\begin{enumerate}
    \item \textbf{Compilazione:} deve compilare senza introdurre
    errori sintattici o di tipizzazione;
    \item \textbf{Correttezza Semantica:} deve risolvere il
    difetto semantico facendo passare i test della
    suite precedentemente falliti;
    \item \textbf{Preservazione Funzionale:} deve preservare integralmente
    il comportamento corretto preesistente del sistema,
    evitando l'introduzione di regressioni.
\end{enumerate}

Per affrontare questa sfida, il \texttt{PatchAgent}
introduce un duplice contributo architetturale:
\begin{itemize}
    \item \textbf{Il ciclo di riparazione iterativo basato su FSM:}
    Il cuore logico dell'agente è governato da
    una Macchina a Stati Finiti (\textit{Finite State
    Machine}, FSM) implicita. Questa struttura dirige
    un agente LLM (dotato di un toolkit
    di sei strumenti specifici) attraverso quattro macro-fasi
    sequenziali: \texttt{EXPLORE}, \texttt{PATCH}, \texttt{COMPILE}
    e \texttt{TEST}. L'agente non si limita a
    un singolo tentativo cieco, ma è dotato
    di un meccanismo di \textit{retry} adattivo,
    il quale sfrutta schedule di temperatura dinamici
    e una memoria conversazionale persistente per imparare
    dai propri fallimenti e aggiustare la strategia
    nei tentativi successivi.
    \item \textbf{La riparazione strutturale basata su Tree-sitter:}
    Abbandonando gli approcci testuali tradizionali (come
    l'uso di espressioni regolari, diff testuali
    o semplici sostituzioni di stringhe), il sistema
    adotta una manipolazione diretta del CST (\textit{Concrete Syntax Tree}) . Questo
    garantisce che ogni singola modifica apportata al
    codice sorgente Java sia corretta dal
    punto di vista posizionale e sintattico per
    costruzione, dimostrandosi robusta indipendentemente dalla complessità
    o dallo stile di codifica del progetto
    target.
\end{itemize}

\subsection{il Ruolo del PatchAgent nell'Architettura HybridRepair}
\label{subsec:3.4.1}
Prima di poter operare, il \texttt{PatchAgent}
necessita di un contesto informativo ricco e strutturato.
Esso riceve in ingresso tre fondamentali artefatti
informativi prodotti dalle fasi precedenti, rigorosamente formalizzati
tramite i dataclass definiti in \texttt{shared/models.py}:

\begin{description}
    \item[\texttt{FaultReport} (dalla Fase 1):] Rappresenta la
    mappa del problema. Contiene le posizioni esatte
    del fault (localizzate dal modulo \texttt{LogicFL}),
    il codice sorgente del metodo difettoso estratto
    in modo sicuro tramite Tree-sitter (da
    \texttt{ast\_extractor.py}), le catene causali ancorate
    semanticamente al codice (da \texttt{prolog\_grounding.py}),
    e infine i test falliti accompagnati dai
    relativi stack trace.
    \item[\texttt{RepairSpec} (dalla Fase 2):] Costituisce il
    piano d'azione. Include una descrizione in linguaggio
    naturale del comportamento difettoso osservato, del
    comportamento corretto atteso e della modifica minimale
    suggerita per la risoluzione, il tutto validato
    criticamente dal modulo \texttt{SpecCritic} che ne
    assegna un punteggio di confidenza.
    \item[\texttt{Ingredients} (dalla Fase 2):] Forniscono il
    materiale da costruzione. Rappresentano l'intero vocabolario
    API reale e accessibile nel contesto della
    classe difettosa, includendo campi, metodi, tipi
    importati e variabili locali, estratti
    dal modulo \texttt{ingredient\_forge.py}.
\end{description}

L'orchestratore principale del sistema, \texttt{repair\_bug.py}, raccoglie
questi tre elementi e istanzia il \texttt{PatchAgent}
nel seguente modo:

\begin{lstlisting}[language=Python , label={lst:init_patchagent}]
agent = PatchAgent(
    fault_report=fault_report,
    repair_spec=repair_spec,
    ingredients=ingredients
)

repair_result = agent.run()
\end{lstlisting}

L'invocazione del metodo principale \texttt{agent.run()} dà
il via al ciclo di riparazione autonomo,
restituendo infine un oggetto \texttt{RepairResult}. Questo
oggetto non contiene solo l'esito finale, ma
aggrega tutti i tentativi effettuati, identifica il
tentativo vincente (qualora la riparazione abbia avuto
successo) e conserva l'intero storico degli artefatti
intermedi, come le conversazioni dell'LLM, i
diff generati e gli output del compilatore.

Il flusso complessivo dell'esecuzione governato dal metodo
\texttt{run()} può essere schematizzato visivamente attraverso
il diagramma a blocchi riportato in Figura~\ref{fig:flusso_patchagent}.

\begin{figure}[H]
\centering
\begin{tikzpicture}[node distance=1.6cm, every node/.style={fill=white, font=\small}, align=center]
    % Definizione stili dei nodi
    \tikzstyle{block} = [draw, rectangle, minimum width=4cm, minimum height=0.8cm, rounded corners=0.1cm, line width=0.6pt]
    \tikzstyle{io} = [draw, trapezium, trapezium left angle=75, trapezium right angle=105, minimum width=4cm, minimum height=0.8cm, line width=0.6pt]
    \tikzstyle{decision} = [draw, diamond, aspect=1.8, minimum width=3cm, minimum height=0.8cm, line width=0.6pt]
    \tikzstyle{arrow} = [thick,->,>=stealth]

    % Posizionamento nodi
    \node (input) [io] {Input: \texttt{FaultReport}, \texttt{RepairSpec}, \texttt{Ingredients}};
    \node (init) [block, below of=input, node distance=1.4cm] {Inizializzazione\\(\texttt{Attempt = 1})};
    \node (setup) [block, below of=init, node distance=1.6cm] {Setup Sandbox \&\\Set Temperature};
    \node (loop) [block, below of=setup, node distance=1.6cm] {\textbf{Tool-Use Loop}\\(LLM + Tools)};
    \node (verify) [block, below of=loop, node distance=1.6cm] {Verifica Finale\\(\texttt{javac} + \texttt{JUnit})};
    \node (dec1) [decision, below of=verify, node distance=1.8cm] {Test Pass?};
    \node (dec2) [decision, left of=dec1, node distance=4.5cm] {Attempt $<$ 3?};
    \node (res) [block, below of=dec1, node distance=2.2cm] {Genera \texttt{RepairResult}};
    \node (output) [io, below of=res, node distance=1.4cm] {Output};

    % Collegamenti frecce
    \draw [arrow] (input) -- (init);
    \draw [arrow] (init) -- (setup);
    \draw [arrow] (setup) -- (loop);
    \draw [arrow] (loop) -- (verify);
    \draw [arrow] (verify) -- (dec1);
    \draw [arrow] (dec1) -- node[anchor=west, xshift=2pt] {YES} (res);
    \draw [arrow] (dec1) -- node[anchor=south, yshift=2pt] {NO} (dec2);
    \draw [arrow] (dec2) |- node[anchor=west, yshift=20pt] {YES} (setup);
    \draw [arrow] (dec2) |- node[anchor=north, xshift=30pt, yshift=-2pt] {NO} (res);
    \draw [arrow] (res) -- (output);

    % Nota sul ciclo FSM affiancata al loop
    \node (fsm_note) [right of=loop, node distance=4.8cm, text width=4cm, font=\footnotesize\itshape, draw, dashed, rounded corners=0.05cm, fill=gray!5] {
        \textbf{Ciclo FSM Coerente:}\\
        \texttt{EXPLORE} $\rightarrow$ \texttt{PATCH} $\rightarrow$ \texttt{COMPILE} $\rightarrow$ \texttt{TEST}
    };
    \draw [thick, dotted, ->] (fsm_note) -- (loop);
\end{tikzpicture}
\caption{Flusso complessivo del ciclo di esecuzione iterativo e adattivo gestito all'interno di \texttt{PatchAgent.run()}.}
\label{fig:flusso_patchagent}
\end{figure}

\subsubsection{Spiegazione del Flusso di Esecuzione}
\label{subsubsec:3.4.1.1}
Il diagramma illustra la natura iterativa e auto-correttiva
del \texttt{PatchAgent}. All'inizio di ognuno dei (massimo) tre tentativi,
il sistema crea una copia isolata del codice difettoso in un sandbox
dedicato (\texttt{pipeline\_results/<bug\_id>/attempt\_k/patched\_source/}), per
evitare che le modifiche di un tentativo inquinino i successivi, e imposta
la temperatura del modello secondo uno schedule adattivo ($0.1$, $0.2$, $0.4$)
che ne aumenta la ``creatività'' a fronte di fallimenti ripetuti.

Preparato l'ambiente, l'esecuzione entra nel \textit{Tool-Use Loop}, dove il
modello e il toolkit interagiscono dinamicamente: è qui che prende vita la FSM,
guidando l'agente a esplorare il codice, generare e applicare una patch e
testarla. L'agente dispone di un budget massimo di round di interazione
(tipicamente 9).\\\\

Usciti dal loop, il sistema procede con una verifica finale deterministica:
invoca il compilatore (\texttt{javac}) e il framework di test (\texttt{JUnit})
sul codice nel sandbox. Se tutti i test passano (\texttt{PASS}), la riparazione
ha successo e il tentativo viene registrato come vincente. Se i test falliscono
o il codice non compila, e restano tentativi disponibili, la storia dei
fallimenti viene salvata nella memoria conversazionale e iniettata nel prompt
successivo per indicare all'LLM cosa evitare; il ciclo riparte con una nuova
copia del codice e una temperatura più elevata. Esauriti i tentativi o raggiunto
il successo, il sistema aggrega le informazioni in un oggetto \texttt{RepairResult}.


\subsection{Il Ciclo di Vita dell'Agente: la FSM di PatchAgent}
\label{subsec:3.4.2}

\label{subsubsec:3.4.2.1}
Il comportamento operativo del \texttt{PatchAgent}, all'interno del suo
ciclo di riparazione, è modellato secondo i principi di una Macchina a
Stati Finiti (FSM, \textit{Finite-State Machine}), formalmente una quintupla
$M = (Q, \Sigma, \delta, q_0, F)$ con $Q$ insieme degli stati, $\Sigma$
alfabeto degli input, $\delta : Q \times \Sigma \to Q$ funzione di transizione,
$q_0$ stato iniziale e $F$ insieme degli stati accettanti.

Calando questo modello nell'architettura del \texttt{PatchAgent}, l'insieme
degli stati operativi è definito come segue:
\[
    Q = \{\texttt{EXPLORE}, \texttt{PATCH}, \texttt{COMPILE}, \texttt{TEST}, \texttt{DONE}, \texttt{RETRY}\}
\]

Il processo ha sempre inizio nello stato di esplorazione,
dunque $q_0 = \texttt{EXPLORE}$, e si conclude con successo
o esaurimento dei tentativi nello stato finale, per cui
$F = \{\texttt{DONE}\}$.
La dinamica di questo automa è descritta dalla funzione
di transizione $\delta$, la quale determina come l'agente avanza
verso la riparazione reagendo all'esito delle sue stesse azioni.

\newpage
La Tabella~\ref{tab:fsm_transitions} illustra le transizioni di stato previste dal sistema:

\begin{table}[H]
\centering
\caption{Transizioni di stato della FSM di \texttt{PatchAgent}.}
\label{tab:fsm_transitions}
\begin{tabular}{p{2.5cm} p{7.5cm} p{3.5cm}}
\toprule
\textbf{Stato \newline corrente} & \textbf{Evento (Input $\Sigma$)} & \textbf{Stato successivo} \\
\midrule
\texttt{EXPLORE} & Strumenti \texttt{read\_file} / \texttt{search\_symbol} / \texttt{get\_method\_signature} eseguiti & \texttt{PATCH} \\
\texttt{PATCH} & Strumento \texttt{apply\_ast\_patch} eseguito con successo & \texttt{COMPILE} \\
\texttt{COMPILE} & \texttt{compile\_check} $\rightarrow$ compilazione riuscita & \texttt{TEST} \\
\texttt{COMPILE} & \texttt{compile\_check} $\rightarrow$ errore di compilazione & \texttt{PATCH} (retry inline) \\
\texttt{TEST} & \texttt{run\_tests} $\rightarrow$ tutti i test passati & \texttt{DONE} \\
\texttt{TEST} & \texttt{run\_tests} $\rightarrow$ test falliti & \texttt{RETRY} (se attempt $<$ 3) \\
\texttt{RETRY} & Inizio nuovo tentativo con temperatura elevata & \texttt{EXPLORE} \\
\texttt{RETRY} & \texttt{attempt} == 3 & \texttt{DONE} (fallimento) \\
\bottomrule
\end{tabular}
\end{table}

\begin{figure}[H]
\centering
\begin{tikzpicture}[->, >=stealth, node distance=2.5cm and 3.5cm, every node/.style={font=\footnotesize},
    state/.style={circle, draw, thick, minimum size=1.8cm, text centered, font=\ttfamily},
    accepting/.style={circle, draw, thick, double, minimum size=1.8cm, text centered, font=\ttfamily}]

    % Nodi
    \node[state] (explore) {EXPLORE};
    \node (start) [left=1.2cm of explore, font=\footnotesize\bfseries] {Inizio};
    \node[state] (patch) [right=of explore] {PATCH};
    \node[state] (compile) [below=of patch] {COMPILE};
    \node[state] (test) [left=of compile] {TEST};
    \node[state] (retry) [left=of test] {RETRY};
    \node[accepting] (done) [below=of test] {DONE};

    % Transizioni
    \draw (start) edge (explore);
    \draw (explore) edge[above] node{Strumenti usati} (patch);
    \draw (patch) edge[bend left=25, right] node{\texttt{apply\_patch}} (compile);
    \draw (compile) edge[bend left=25, left] node{Errore comp.} (patch);
    \draw (compile) edge[above] node{Compilazione ok} (test);
    \draw (test) edge[right] node{Test \texttt{PASS}} (done);
    \draw (test) edge[above] node{Test \texttt{FAIL} ($<$ 3)} (retry);
    \draw (retry) edge[bend left=30, left=5pt] node{Nuovo tentativo} (explore);
    \draw (retry) |- node[near start, left] {Test \texttt{FAIL} ($= 3$)} (done);
\end{tikzpicture}
\caption{Grafo delle transizioni della Macchina a Stati Finiti (FSM) implicita del \texttt{PatchAgent}.}
\label{fig:fsm_grafo}
\end{figure}

\newpage
\paragraph{Il Concetto di FSM Emergente}
Un aspetto chiave del design di \texttt{HybridRepair} è che questa FSM non è
implementata in modo tradizionale: nel codice non esiste una classe Python con
un attributo \texttt{state} esplicito né una tabella di transizioni codificata
con istruzioni \texttt{if/switch}. Essa prende vita organicamente dalla combinazione del \textit{System Prompt}
iniettato nel Large Language Model (GPT-4o) e dal loop
di controllo implementato nei metodi \texttt{PatchAgent.run()} e \texttt{AzureClient.chat\_with\_tools()}. Sulla
scorta dei successi ottenuti da sistemi recenti come \texttt{RepairAgent}
e \texttt{ChatRepair}, che hanno dimostrato la sorprendente capacità dei
modelli GPT-4 di aderire a \textit{workflow} multi-step complessi tramite
\textit{prompting} strutturato, il modello di \texttt{HybridRepair} impara a transitare
autonomamente da uno stato all'altro, selezionando dinamicamente lo strumento
più appropriato in ciascun round conversazionale \cite{xia2023chatrepair, bouzenia2025repairagent}.

Questa coreografia implicita è diretta da un prompt di
sistema, il cui nucleo normativo è riportato di seguito:

\begin{lstlisting}[language=Python, label={lst:agent_system_prompt}]
AGENT_SYSTEM_PROMPT = """\
You are an expert Java bug repair agent. Your task is to fix a bug in a \
Defects4J Java project using the tools available to you.

## Workflow
1. **EXPLORE**: Read the buggy code, search for related symbols, check method signatures.
2. **PATCH**: Apply a fix using `apply_ast_patch`. Provide the COMPLETE corrected method.
3. **COMPILE**: Check compilation with `compile_check`.
4. **TEST**: Verify with `run_tests`.

## Rules
- ALWAYS use `apply_ast_patch` to apply fixes. Never output raw diffs.
- ALWAYS compile before testing.
- Use REAL API methods only. Call `get_method_signature` to verify before using an API.
- If compilation fails, read the error, fix the issue, and try again.
- If tests fail, analyze the stack trace and try a different approach.
- Keep fixes MINIMAL. Prefer the smallest change that fixes the bug.
- Do NOT add unnecessary null checks or try/catch blocks.
...
"""
\end{lstlisting}

\newpage
\paragraph{Analisi delle Regole Fondamentali del Prompt}
Ogni singolo elemento di questo \texttt{AGENT\_SYSTEM\_PROMPT} è stato concepito
per vincolare semanticamente l'output generativo del modello, trasformando un
LLM \textit{general-purpose} in un risolutore di bug. Vi
sono tre regole fondamentali che meritano un'analisi approfondita:

\begin{description}
    \item[La Sequenza Numerata delle Azioni:]
    (\texttt{1.\ EXPLORE} $\rightarrow$ \texttt{2.\ PATCH} $\rightarrow$ \texttt{3.\ COMPILE} $\rightarrow$ \texttt{4.\ TEST}):
    Questa enumerazione codifica il grafo di transizione della FSM in linguaggio naturale.
     I modelli Transformer (come \texttt{GPT-4o}) sono intrinsecamente \textit{stateless} e non
     possiedono una rappresentazione nativa di un automa a stati finiti \cite{sumers2023cognitive};
     tuttavia, nei dati di addestramento una lista numerata è fortemente associata a una sequenza
     causale e temporale di azioni. Sfruttando questa proprietà \cite{wang2023plan}, il modello
     è indotto a procedere in modo ordinato, esplorando l'ambiente prima di applicare una modifica,
     in linea con i paradigmi di \textit{Reasoning and Acting} \cite{yao2022react}.

    \item[Il Divieto di Diff Testuali:]
    (\texttt{"ALWAYS use apply\_ast\_patch... Never output raw diffs."}):
    Se non vincolati, i modelli tendono a eludere l'invocazione degli strumenti, ripiegando su
    frammenti di codice o diff testuali inseriti nella prosa. La generazione diretta di \textit{raw diffs}
    è però prona a errori di formattazione e incompatibile con le pipeline di applicazione
    automatizzata \cite{jimenez2024sweagent}. La negazione esplicita (\textit{negative constraint})
    forza l'agente a delegare la manipolazione del codice unicamente a \texttt{apply\_ast\_patch}
    \cite{qin2023toolllm}: senza questo vincolo il modello potrebbe descrivere a parole una soluzione
    corretta senza applicarla, vanificando la fase di \texttt{COMPILE} poiché il sorgente resterebbe inalterato.

    \item[Il Vincolo di Minimalità:]
    (\texttt{"Keep fixes MINIMAL..."}):
    Questa direttiva recepisce un principio cardine dell'APR \cite{huang2024evolving}: le patch minime
    hanno una probabilità inferiore di indurre regressioni su comportamenti non coperti dalla suite di
    test. Ricerche su sistemi allo stato dell'arte \cite{xia2024fitrepair} mostrano che le patch
    \textit{single-hunk} (concentrate su un singolo blocco contiguo) sono più plausibili e generalizzabili
    delle modifiche \textit{multi-hunk} distribuite. Privilegiare la minima perturbazione necessaria
    mitiga l'overfitting sui test falliti \cite{campos2025exploring}.
\end{description}


\subsubsection{La Fase EXPLORE: Esplorazione Autonoma del Workspace}
\label{subsubsec:3.4.2.2}
La prima macro-fase in cui la Macchina a Stati Finiti
fa il suo ingresso è la fase \texttt{EXPLORE}. In questo frangente,
l'agente LLM ha il compito di costruire una solida comprensione
contestuale del difetto prima di tentare di alterare il codice.
Durante questa fase esplorativa, l'agente è libero di invocare tre
strumenti specifici del suo toolkit: \texttt{read\_file}, \texttt{search\_symbol} e \texttt{get\_method\_signature}.

La fase di esplorazione potrebbe sembrare ridondante: il prompt iniziale fornisce già il codice
del metodo difettoso, la \texttt{RepairSpec} e il vocabolario delle API (\texttt{Ingredients}).
Tuttavia queste informazioni sono localizzate attorno all'epicentro del difetto, e due categorie
di anomalie rendono insufficiente questo perimetro:

\begin{description}
    \item[Bug con propagazione cross-method:] La causa radice non sempre risiede nel metodo in cui
    l'eccezione viene sollevata: può nascondersi in un metodo chiamato a cascata o in interfacce
    implementate da una sottoclasse. In questi casi l'agente deve navigare oltre il file iniziale,
    leggendo moduli aggiuntivi per mappare il flusso di esecuzione prima di formulare una patch
    semanticamente risolutiva.

    \item[Bug con API ambigue:] Anche quando la \texttt{RepairSpec} suggerisce una specifica API, la
    firma esatta può non essere chiara a priori (overload multipli, parametri generici con vincoli).
    Lo strumento \texttt{get\_method\_signature} permette all'agente di ispezionare le firme reali e
    disambiguare prima di sintetizzare il codice.
\end{description}

La durata della fase \texttt{EXPLORE}, misurata in numero di round di invocazione
degli strumenti, non è prefissata: l'agente valuta autonomamente quando ha
raccolto informazioni sufficienti per agire. La transizione a \texttt{PATCH} è
implicita, innescata nel momento in cui l'output del modello include
un'invocazione allo strumento \texttt{apply\_ast\_patch}.

\subsubsection{La Fase PATCH: Applicazione Chirurgica della Modifica}
\label{subsubsec:3.4.2.3}
Nella fase \texttt{PATCH} l'agente invoca lo strumento \texttt{apply\_ast\_patch},
fornendo come argomenti il percorso del file bersaglio, il nome del metodo da
rimpiazzare, il codice sorgente completo del nuovo metodo e, se necessario, un
numero di riga come \textit{hint} per risolvere ambiguità di overloading.

In questa fase vige il vincolo più rigoroso del sistema: il modello deve fornire il metodo nella
sua interezza (firma, graffe e corpo). Diff unificati o frammenti parziali sono vietati, come ribadito
nel prompt di sistema (\texttt{"Provide the COMPLETE corrected method"}) e nella definizione dello strumento.
Questa scelta, all'apparenza dispendiosa in token, è sorretta da tre motivazioni:

\begin{description}
    \item[Completezza sintattica garantita:] I diff testuali sono fragili quando pilotati da LLM: i
    modelli faticano a generare righe di contesto esatte e a rispettare l'indentazione, incorrendo in
    allucinazioni strutturali \cite{zhang2024llmhallucinations}, e in un ciclo iterativo un contesto
    imperfetto fa fallire subito l'applicazione. Riscrivere il metodo per intero sposta il problema dalla
    manipolazione testuale alla sostituzione semantica a livello di AST \cite{pabba2025semagent}, eliminando
    le ambiguità di conteggio righe e formattazione, con un tasso di successo superiore ai diff riga per
    riga \cite{xia2023chatrepair}.

    \item[Compatibilità con Tree-sitter:] L'infrastruttura di manipolazione (\texttt{ast\_applier.py}) opera
    sul \textit{Concrete Syntax Tree} (CST) di Tree-sitter, individuando il nodo del metodo (\texttt{method\_declaration}
    o \texttt{constructor\_declaration}) e isolandone il \textit{text span}. Per sostituire quel text span l'input
    deve essere un metodo Java completo e autonomo: frammenti isolati romperebbero l'albero.

   \item[Prevenzione di patch parziali:] Gli LLM tendono a risposte ``pigre'' o iper-focalizzate
    \cite{zhang2024llmhallucinations}: istruiti a generare solo la modifica, si concentrano sulla riga
    dell'errore (es. un \texttt{if (obj != null)}) ma trascurano il bilanciamento delle graffe o la
    rilocazione di variabili dichiarate altrove. Imporre la generazione dell'intero metodo mantiene la
    struttura logica nell'orizzonte di attenzione del modello, riducendo errori di scope e sintassi
    sbilanciata \cite{xia2023chatrepair}.
\end{description}

\subsubsection{La Fase COMPILE: Feedback Deterministico dal Compilatore}
\label{subsubsec:3.4.2.4}
Una volta che la patch è stata applicata chirurgicamente
al codice sorgente tramite l'AST, l'agente transita nella fase
\texttt{COMPILE}, che rappresenta il primo vero ``oracolo deterministico''
dell'intero ciclo di riparazione.

In questa fase emerge la differenza tra l'agente generativo e gli strumenti di
validazione. Il giudizio dell'LLM sulla correttezza del codice è probabilistico
e vulnerabile ad allucinazioni e sviste sintattiche; il compilatore \texttt{javac},
al contrario, applica in modo esatto la grammatica del linguaggio e dà un verdetto
binario: la patch è valida sotto il profilo della tipizzazione statica, oppure
viene rifiutata.

L'interazione tra l'agente e questo oracolo avviene invocando lo
strumento \texttt{compile\_check}. Sebbene l'implementazione completa si appoggi al modulo
\texttt{sandbox\_evaluator.py}, la logica di interfaccia esposta all'LLM può essere
riassunta nel seguente frammento di codice:

\begin{lstlisting}[language=Python, label={lst:tool_compile_check}]
def tool_compile_check(context: ToolContext) -> str:
    """Verifica la compilazione del file sorgente modificato."""
    result = compile_patched_source(context.patched_dir, context.bug_dir)
    
    if result.is_success:
        return "COMPILATION SUCCESS"
    else:
        # Troncamento diagnostico per preservare il budget di token
        error_msg = result.stderr[:3000] 
        return f"COMPILATION ERROR:\n{error_msg}"
\end{lstlisting}

Un aspetto cruciale visibile in questa funzione è
il troncamento dell'output a 3000 caratteri. Quando \texttt{javac} incontra
un errore sintattico in un file Java (specialmente se
l'errore risiede in firme di metodi o parentesi non
chiuse correttamente), tende a produrre una cascata di errori
secondari derivati estremamente verbosa. Inviare decine di migliaia di
caratteri di log non solo consumerebbe rapidamente la finestra
di contesto del modello (\textit{token budget}), ma aggiungerebbe ``rumore''
non necessario. Troncando l'output, il sistema si assicura di
fornire all'agente solo i primissimi errori (le cause radice),
presentati nel noto formato standard \texttt{NomeFile.java:LineaNum: error: <tipo\_errore> <descrizione>}.
Il modello GPT-4o, ampiamente addestrato su log di compilazione,
è in grado di estrarre autonomamente il numero di
riga da questa stringa per risalire al punto esatto
del fallimento.

Questo feedback innesca il ciclo di \textit{retry inline} (\texttt{PATCH} $\rightarrow$ \texttt{COMPILE} $\rightarrow$ \texttt{PATCH}):
se la compilazione fallisce, il messaggio di errore viene reiniettato nel contesto e l'agente corregge la
patch nello stesso round con una nuova invocazione ad \texttt{apply\_ast\_patch}. Il micro-ciclo consuma il
budget di round (\texttt{max\_tool\_rounds}) ma permette di correggere sviste sintattiche minori senza
scartare l'intera strategia.

\subsubsection{La Fase TEST: Verifica Semantica via Suite JUnit}
\label{subsubsec:3.4.2.5}
Se la fase \texttt{COMPILE} viene superata con successo, la
FSM avanza verso il suo stato più critico: la
fase \texttt{TEST}. Se il compilatore garantisce che il codice
può essere eseguito senza violare le regole del linguaggio,
il framework di test JUnit funge da secondo oracolo,
certificando la correttezza semantica e comportamentale della patch. In
altre parole, JUnit verifica se il programma modificato risolve
effettivamente il difetto senza alterare le funzionalità preesistenti.

L'agente innesca questa verifica tramite lo strumento \texttt{run\_tests}, la
cui logica si appoggia anch'essa al sandbox di valutazione:

\begin{lstlisting}[language=Python, label={lst:tool_run_tests}]
def tool_run_tests(context: ToolContext) -> str:
    """Esegue la suite JUnit per validare la semantica della patch."""
    test_result = evaluate(context.patched_dir, context.test_classes)
    
    if test_result.status == "PASS":
        return "PASS: All tests passed successfully."
        
    elif test_result.status == "COMPILE_ERROR":
        return f"COMPILE_ERROR:\n{test_result.error_output}"
        
    elif test_result.status == "TEST_FAIL":
        return (f"TEST_FAIL:\n{test_result.stack_traces}\n"
                f"Previously failing now passing: {test_result.previously_failing_now_passing}")
\end{lstlisting}

L'esecuzione dei test può condurre a tre scenari di
output ben distinti, che forniscono all'agente un feedback semanticamente
ricco:

\begin{description}
    \item[\texttt{PASS}:] Rappresenta il successo totale. Tutti i test della
    suite, inclusi quelli che originariamente evidenziavano il bug, sono
    stati eseguiti senza sollevare eccezioni. In questo scenario, la
    FSM transita direttamente allo stato \texttt{DONE} e la riparazione
    è considerata vincente.
    
    \item[\texttt{COMPILE\_ERROR}:] Uno scenario atipico ma possibile, che si verifica
    se la fase di verifica finale (che ricompila il
    progetto prima di lanciare JUnit) individua difetti sintattici residui
    o disallineamenti non intercettati dal precedente \texttt{compile\_check} isolato. Comunica
    all'agente che il codice, allo stato attuale, non è
    eseguibile.
    
    \item[\texttt{TEST\_FAIL}:] Indica che la patch è compilabile (valida sintatticamente),
    ma semanticamente scorretta o incompleta. In questo caso, lo
    strumento restituisce all'LLM l'intero stack trace delle eccezioni sollevate
    dai test falliti. L'agente può studiare questi stack trace
    per comprendere per quale input o in quale condizione
    logica la sua patch ha ceduto, permettendogli di riconsiderare
    il proprio approccio.
\end{description}

Di particolare importanza metrica all'interno dello scenario di fallimento
è il calcolo del campo \texttt{previously\_failing\_now\_passing}. Estratto e calcolato
dalla funzione \texttt{\_parse\_junit\_output()}, questo valore indica quanti e quali test,
tra quelli che nel \texttt{FaultReport} iniziale fallivano inesorabilmente, sono
stati portati al successo dalla patch corrente. Questa informazione
è vitale per misurare i progressi parziali: una correzione
che riesce a far passare 3 test su 5
precedentemente falliti rappresenta un chiaro indicatore direzionale positivo. Anche
se non soddisfa il criterio di successo assoluto, dimostra
che la strategia adottata sta parzialmente risolvendo il problema,
fornendo preziose indicazioni per i tentativi futuri.

\subsubsection{Il Meccanismo di Retry: Schedule di Temperatura e Adattamento}
\label{subsubsec:3.4.2.6}
Quando la fase \texttt{TEST} fallisce, il ciclo non si interrompe: la FSM entra
nella fase transizionale di \texttt{RETRY}, che prepara un nuovo tentativo. Per
evitare di ricadere negli stessi errori, il meccanismo di retry si fonda su tre
pilastri: una modulazione adattiva della temperatura del modello, l'isolamento
del workspace e l'iniezione della memoria conversazionale dei fallimenti pregressi.

\paragraph{1. Schedule di Temperatura Adattiva}
Il primo pilastro governa la parametrizzazione del \textit{Large Language Model}.
La temperatura di decodifica, il parametro stocastico che
controlla l'appiattimento della distribuzione di probabilità sui token in
fase di campionamento, non è statica, ma segue
uno schedule progressivo codificato nel modulo \texttt{shared/config.py}:

\begin{lstlisting}[language=Python, label={lst:temp_schedule}]
# shared/config.py
LLM_TEMPERATURE_SCHEDULE = [0.1, 0.2, 0.4]
COMPILE_ERROR_TEMP_DOWNSHIFT = 0.1
\end{lstlisting}

La scelta di questi valori specifici (\texttt{0.1}, \texttt{0.2} e
\texttt{0.4}) risponde a una logica di esplorazione calibrata:

\begin{description}
    \item[\texttt{0.1} (Primo tentativo):] Una temperatura bassa forza un regime deterministico, che
    favorisce i token più probabili; per bug ben diagnosticati dalla Fase 1, la soluzione più ovvia è
    spesso quella corretta.

    \item[\texttt{0.2} (Secondo tentativo):] Dopo un primo fallimento, l'innalzamento introduce una
    marginale diversità semantica per esplorare approcci alternativi, restando conservativo.

    \item[\texttt{0.4} (Terzo tentativo):] Nel tentativo finale abilita una diversificazione marcata, utile
    a sfuggire ai \textit{local minima} (dove il modello ripropone varianti dello stesso codice errato),
    pur preservando la correttezza sintattica.
\end{description}

In aggiunta, il parametro \texttt{COMPILE\_ERROR\_TEMP\_DOWNSHIFT} introduce un'ottimizzazione adattiva reattiva:
se il tentativo precedente è fallito non a livello
semantico (fase \texttt{TEST}) ma sintattico (fase \texttt{COMPILE}), significa che il
modello si è già allontanato troppo dalla grammatica valida.
In questo scenario, il sistema decurta preventivamente la temperatura
per il tentativo successivo, forzando una forte convergenza verso
costrutti sintattici più sicuri e ortodossi.

\paragraph{2. Isolamento della Sorgente per Ogni Tentativo}
Il secondo pilastro garantisce la riproducibilità e la pulizia
del contesto di lavoro. Per impedire che l'inquinamento del
codice sorgente causi falsi negativi, il sistema esegue un
isolamento fisico su disco prima di instradare il nuovo
tentativo:

\begin{lstlisting}[language=Python, label={lst:create_patched_copy}]
def _create_patched_copy(self, attempt_num: int) -> Path:
    """
    Crea una directory isolata e una copia pulita del codice sorgente
    per il tentativo corrente.
    """
    patched_dir = self.results_dir / f"attempt_{attempt_num}" / "patched_source"
    if patched_dir.exists():
        shutil.rmtree(patched_dir)
    
    # Copia profonda per isolare completamente il tentativo
    shutil.copytree(self.bug_dir, patched_dir)
    return patched_dir
\end{lstlisting}

\texttt{shutil.copytree()} realizza una copia profonda e ricorsiva del progetto target, garantendo che
ogni tentativo parta dal file difettoso originale, senza che patch fallite di tentativi precedenti
contaminino le iterazioni successive. Conservare in directory separate (\texttt{attempt\_1}, \texttt{attempt\_2},
ecc.) lo stato di ogni iterazione fornisce inoltre un artefatto utile all'ispezione \textit{post-hoc}.

\paragraph{3. Copia Impostata degli Errori Precedenti}
Il terzo pilastro è rappresentato dalla memoria a breve
termine, essenziale per abilitare il meta-apprendimento dell'agente, che viene
delegata a un modulo specifico.
\newpage
\subsubsection{La Memoria Conversazionale: \texttt{conversation\_store.py}}
\label{subsubsec:3.4.2.7}
I sistemi APR multi-tentativo di prima generazione soffrivano di un limite:
l'assenza di ritenzione dello storico tra iterazioni. Ogni tentativo ricominciava
dal medesimo prompt iniziale, senza comunicare quali patch fossero già state
tentate né perché fossero fallite. Questo approccio ``smemorato'' produceva il
cosiddetto \textit{loop sterile}: in assenza di variazioni sostanziali del prompt,
il modello rigenerava la stessa patch a ogni tentativo, consumando il budget
computazionale senza progressi.

Per scardinare questo loop sterile, \texttt{HybridRepair} adotta ed estende
il paradigma presente in \textit{ChatRepair}\cite{xia2023chatrepair}. Questo schema prevede che la storia
conversazionale associata a ogni tentativo fallito venga persistita e
successivamente distillata in un riassunto iniettato nel prompt del
tentativo seguente come contesto negativo. Il modulo \texttt{patch\_agent/conversation\_store.py} funge
da database locale per questa memoria a breve termine.

Per massimizzare la resilienza I/O, la storicizzazione della conversazione
avviene sfruttando il formato JSONL (\textit{JSON Lines}), implementato attraverso
la funzione \texttt{save()}:

\begin{lstlisting}[language=Python, label={lst:conversation_save}]
def save(self, attempt_num: int, messages: list):
    """
    Serializza i messaggi del tentativo in formato JSONL
    per garantire tolleranza ai guasti durante l'I/O.
    """
    file_path = self.store_dir / f"attempt_{attempt_num}.jsonl"
    with open(file_path, "w", encoding="utf-8") as f:
        for msg in messages:
            f.write(json.dumps(msg) + "\n")
\end{lstlisting}

La scelta del JSONL al posto di un array JSON monolitico è strutturale: poiché ogni messaggio (ruolo,
contenuto e chiamate agli strumenti) è un record isolato su una riga, un'interruzione del processo non
corrompe i messaggi già scritti, mentre in un array JSON la parentesi di chiusura mancante invaliderebbe
l'intero file.

Quando l'agente sta per iniziare il tentativo $k+1$, il
sistema interroga il \texttt{conversation\_store.py} per costruire il ponte cognitivo
tra il fallimento passato e il nuovo sforzo, invocando
la funzione \texttt{get\_failure\_summary()}:

\begin{lstlisting}[language=Python, label={lst:get_failure_summary}]
def get_failure_summary(self) -> str:
    """
    Costruisce un riassunto dei tentativi falliti passati
    utilizzando il costrutto di negative prompting.
    """
    failed_attempts = self._load_previous_failures()
    if not failed_attempts:
        return ""

    summary = "DO NOT REPEAT THESE APPROACHES:\n\n"
    for attempt in failed_attempts:
        summary += f"Attempt {attempt.id} Action:\n"
        summary += f"{attempt.last_assistant_message}\n"
        summary += f"Resulting Error:\n{attempt.latest_error}\n\n"
        
    return summary
\end{lstlisting}

Invece di chiedere genericamente di ``cercare un'altra via'', la stringa esplicita quale azione ha
compiuto l'assistente nell'ultimo messaggio del tentativo fallito e quale errore ha scatenato.
La testata \texttt{"DO NOT REPEAT THESE APPROACHES"} agisce da \textit{negative prompting} forte: fornire
esempi espliciti con un imperativo negativo penalizza le sequenze simili al testo negato, spingendo
l'agente verso ragionamenti inesplorati e rendendo le iterazioni effettivamente progressive.

